\setcounter{section}{0}

\Needspace{5\baselineskip}
\hypertarget{ux81eaux56deux5f52ux6a21ux578b}{%
\section{自回归模型}\label{ux81eaux56deux5f52ux6a21ux578b}}

\Needspace{5\baselineskip}
\hypertarget{ux4e00ux6838ux5fc3ux601dux60f3}{%
\subsection{一、核心思想}\label{ux4e00ux6838ux5fc3ux601dux60f3}}

一个序列中当前时刻的值，可以由其过去若干个时刻的值的线性组合加上一个随机扰动（误差）来解释和预测。如基于这个上下文来预测最可能出现的``下一个''词。

\Needspace{5\baselineskip}
一个~p阶~的自回归模型，记为~AR(p)，其公式定义为：

\[
X_t = c + \phi_1 X_{t-1} + \phi_2 X_{t-2} + \cdots + \phi_p X_{t-p} + \epsilon_t
\]

\Needspace{5\baselineskip}
公式解读：

\begin{itemize}
\tightlist
\item
  \(X_t\)：当前时刻 \(t\) 的变量值（我们想要预测的值）。
\item
  \(c\)：常数项（截距）。
\item
  \(\phi_1, \phi_2, \ldots, \phi_p\)：自回归系数。这些是模型需要学习的参数，代表了过去每个时刻的值对当前值的影响程度和方向。
\item
  \(X_{t-1}, X_{t-2}, \ldots, X_{t-p}\)：过去 \(p\) 个时刻的历史值。
\item
  \(\epsilon_t\)：时刻 \(t\) 的随机误差项（白噪声）。它代表了模型无法用历史数据解释的随机波动，通常假设其均值为 0，方差恒定。
\end{itemize}

阶数 p决定了模型回看多远的过去。p 的选择至关重要，太小可能欠拟合，太大可能过拟合。

平稳性要求：序列的均值、方差和协方差不随时间推移而发生显著变化。如果不平稳，通常需要通过差分等方法将其转换为平稳序列（这就引出了ARIMA模型中的``I''）。

\Needspace{5\baselineskip}
\hypertarget{ux4e8cux5e94ux7528}{%
\subsection{二、应用}\label{ux4e8cux5e94ux7528}}

\begin{enumerate}
\def\labelenumi{\arabic{enumi}.}
\tightlist
\item
  LLM：将生成问题看作一个序列生成问题（词序列、token序列）。模型定义了一个条件概率分布：\(P(x_t | x_{<t})\)。即，给定所有已经生成的词 \(x_1, x_2, ..., x_{t-1}\)，计算下一个词 \(x_t\) 的概率分布。从该分布中采样（或取最可能）得到下一个词 \(x_t\)。将新生成的词 \(x_t\) 加入上下文，重复步骤，直到生成序列结束符或达到最大长度。
\end{enumerate}

实现技术：~通过掩码机制（Causal Masking）~实现。在Transformer中，确保在计算第t个位置的注意力时，只能``看到''第1到第t-1个位置的信息，无法看到``未来''的信息，从而严格保证了自回归的特性。

\begin{enumerate}
\def\labelenumi{\arabic{enumi}.}
\setcounter{enumi}{1}
\tightlist
\item
  计算机视觉：也可以用``给定之前所有生成的概率，预测下一个像素的概率分布''的方法生成非常清晰、高质量的图像，但缺点是非常慢，因为必须逐个像素顺序生成，无法并行。
\end{enumerate}

\FloatBarrier
\Needspace{5\baselineskip}
\hypertarget{ux53c2ux8003ux6587ux732e-17}{%
\subsection{参考文献}\label{ux53c2ux8003ux6587ux732e-17}}

\vspace{0.25em}
\begin{itemize}
\tightlist
\item
  Zhang, A., Lipton, Z. C., Li, M., \& Smola, A. J. (2023). \href{https://D2L.ai}{Dive into Deep Learning}. Cambridge University Press.
\end{itemize}

\clearpage
